\scanpage{25}
\mspage{\unsure{10|19}}
\deleted{\unsure{de $R$}. En vertu de 10.5.3. c'est un sous-groupe fermé de $G$ lisse sur $S$, donc \unsure{l'assertion affirmée}, et comme $T$ est central, on a \unsure{$R\subset H$}.}

Ceci résulte du \deleted{résultat suivant}

\uline{Lemme} 1.14. Soient $G$ \deleted{affine}\inserted{un $S$-préschéma en groupes,} de présentation finie sur $S$, $T$ un tore maximal de $G$, \deleted{alors $T$ contient tout sous-tore central} $R$ un sous-tore de $G$, et supposons que $T$ et $R$ commutent. Alors $R\subset T$.

En effet, comme $R$ et $T$ commutent, le morphisme $R\times_S T\longrightarrow G$ défini \deleted{\unsure{par}} par $(r,t)\rightsquigarrow r.t$ est un homomorphisme de groupes, donc en vertu de 1.8. il admet un sous-groupe image $T'$ dans $G$, qui est un \deleted{sous-}groupe de type multiplicatif quotient de $R\times_S T$, donc un tore, contenant évidemment $T$. Comme $T$ est un tore maximal, on a $T'=T$ (1.4.), donc $R\subset T$, \deleted{cqfd} ce qui démontre 1.14. donc 1.13. En particulier, \inserted{utilisant 6.7. b):}

\uline{Corollaire} 1.15. Soit $G$ un $S$-préschéma en groupes \deleted{\unsure{affi}} \uline{commutatif}, lisse et affine sur $S$, \inserted{de rang réductif localement constant.} Alors $G$ admet un unique tore maximal, et ce\deleted{\unsure{ci}} \deleted{\unsure{sous-tore de plus}} dernier contient tout sous-tore de $G$.

\uline{Corollaire} 1.16. Soit $G$ un $S$-préschéma en groupes, lisse et affine sur $S$. Supposons que pour tout $s\in S$, désignant par $\bar s$ le spectre d'une clôture algébrique \inserted{$k$} de $k(s)$, la fibre géométrique $G_{\bar s}$ soit un groupe algébrique connexe, et nilpotent (au sens de BIBLE, i.e. le groupe $G(k)$ de ses points à valeurs dans $k$ est nilpotent). \inserted{Supposons de plus le rang réductif de $G$ loc.constant.} Alors $G$ admet un unique tore maximal $T$, \inserted{de plus} $T$ est central et c'est le plus grand sous-tore de $G$.

En vertu de 1.7. b), $G$ admet localement pour fpqc un tore maximal, et en vertu de 1.13. on est ramené à prouver qu'un tore maximal de $G$ est central. En vertu de Exp \inserted{$\overline{\underline{\mathrm{IX}}}$}\deleted{\unsure{XI}}, 5.6. b), on est ramené au cas où $S$ est le spectre d'un corps, qu'on peut supposer algébriquement clos.
\mspage{24}
